\section{Discussion}
\label{sec:discussion}

\subsection{Why the Proposed Framework Works}
\label{sec:disc_framework}
Phase~1 decouples descriptor tuning from classifier choice, avoiding co-optimisation bias.
The two-phase design makes feature quality measurable before committing to expensive hyperparameter sweeps across eight model families.
\placeholder{Expand framework rationale.}

\subsection{Impact of Separability-Guided Optimization}
\label{sec:disc_separability}
High composite scores for LBP and HOG align with strong Phase~2 performance when those blocks dominate the feature vector.
Agreement between Phase~1 rankings and Phase~2 importance analysis (Section~\ref{sec:interp_link}) validates the separability-first protocol.
\placeholder{Expand impact analysis.}

\subsection{Comparison with Literature}
\label{sec:disc_literature}
KNN on Full(opt) (F1 $=0.946$) and SVM on Full(opt) (F1 $=0.912$) demonstrate that separability-tuned radiomics on $\sim3{,}000$ imbalanced slices can match many published deep-learning benchmarks on the same public dataset~\cite{kumar2024,ottoni2025,basthikodi2024}.
\placeholder{Add comparison table against recent handcrafted and hybrid MRI studies.}

\subsection{Practical Implications}
\label{sec:disc_practical}
The proposed pipeline offers:
\begin{itemize}
    \item low computational requirements (no GPU training);
    \item explainable feature engineering via named descriptor families;
    \item suitability for limited-data and single-centre environments;
    \item reproducible benchmarking with statistical significance testing.
\end{itemize}
\textbf{Practical recommendation:} begin with separability-guided texture extraction and a kernel SVM or KNN on HOG or Full(opt) features when labelled data are limited.
